\documentclass[12pt,a4paper]{article}

\usepackage[romanian]{babel}
\usepackage[T1]{fontenc}
\usepackage[utf8]{inputenc}
\usepackage{lmodern}
\usepackage{enumitem}

\title{Normalizarea Bazelor de Date \\ Formele Normale 1NF--6NF}
\author{}
\date{}

\begin{document}

\maketitle

\section{Normalizarea}
Normalizarea este un proces de organizare a datelor într-o bază de date relațională
cu scopul de a elimina redundanța, anomaliile de actualizare și inconsistențele.

\section{Prima Formă Normală (1NF)}
O relație se află în \textbf{1NF} dacă:
\begin{itemize}[noitemsep]
    \item toate atributele conțin valori atomice;
    \item fiecare rând este unic;
    \item nu există atribute multivaloare sau compuse.
\end{itemize}

\section{A Doua Formă Normală (2NF)}
O relație este în \textbf{2NF} dacă:
\begin{itemize}[noitemsep]
    \item este în 1NF;
    \item toate atributele non-cheie depind complet de cheia primară;
    \item nu există dependențe parțiale.
\end{itemize}

\section{A Treia Formă Normală (3NF)}
O relație este în \textbf{3NF} dacă:
\begin{itemize}[noitemsep]
    \item este în 2NF;
    \item nu există dependențe tranzitive între atributele non-cheie.
\end{itemize}

\section{Superchei, Chei Candidate și Cheia Primară}

\subsection{Supercheie}
O supercheie este orice set de atribute care identifică unic fiecare tuplu dintr-o relație.

\subsection{Chei Candidate}
Cheile candidate sunt superchei minimale, din care nu se poate elimina niciun atribut.

\subsection{Cheia Primară}
Cheia primară este cheia candidată aleasă pentru identificarea oficială a tuplurilor.

\section{Forma Normală Boyce--Codd (BCNF)}
O relație este în BCNF dacă pentru orice dependență funcțională
\( X \rightarrow Y \), X este supercheie.

\section{A Patra Formă Normală (4NF)}
O relație este în 4NF dacă:
\begin{itemize}[noitemsep]
    \item este în BCNF;
    \item nu conține dependențe multivaloare independente.
\end{itemize}

\section{A Cincea Formă Normală (5NF)}
O relație este în 5NF dacă:
\begin{itemize}[noitemsep]
    \item este în 4NF;
    \item nu poate fi descompusă în tabele mai mici fără pierdere de informație.
\end{itemize}

\section{A Șasea Formă Normală (6NF)}
O relație este în 6NF dacă:
\begin{itemize}[noitemsep]
    \item nu mai poate fi descompusă deloc fără pierdere de informație;
    \item fiecare tabel conține o singură dependență.
\end{itemize}

\end{document}
